\documentclass[aspectratio=169]{slides}
\begin{document}

\begin{frame}{Centred by default}
A short body sits midway between the title and the bottom margin.
\end{frame}

\begin{frame}[t]{Top}
This body hugs its title, as a dense reference frame wants.
\end{frame}

\begin{frame}[b]{Bottom}
This body sits on the bottom margin, leaving the air above.
\end{frame}

\begin{frame}[t]
A bare top frame keeps its first line where an article would put it.
\end{frame}

\begin{frame}{Overflow}
A frame taller than its stage takes nothing above the content: it stays
top flush and spills to a continuation page, never off the top edge.
The placeholder committee resolved to enumerate the resolutions again.
The placeholder committee resolved to enumerate the resolutions again.
The placeholder committee resolved to enumerate the resolutions again.
The placeholder committee resolved to enumerate the resolutions again.
The placeholder committee resolved to enumerate the resolutions again.
The placeholder committee resolved to enumerate the resolutions again.
The placeholder committee resolved to enumerate the resolutions again.
The placeholder committee resolved to enumerate the resolutions again.
The placeholder committee resolved to enumerate the resolutions again.
The placeholder committee resolved to enumerate the resolutions again.
The placeholder committee resolved to enumerate the resolutions again.
The placeholder committee resolved to enumerate the resolutions again.
The placeholder committee resolved to enumerate the resolutions again.
The placeholder committee resolved to enumerate the resolutions again.
The placeholder committee resolved to enumerate the resolutions again.
The placeholder committee resolved to enumerate the resolutions again.
\end{frame}

\end{document}
